\section*{الفصل \ref{ch:b3:inorganic-materials} --- المواد غير العضوية والمواد النانوية}

\begin{solution}{exo:b3:inorganic-materials:1}
$x^2 \propto t$: فسمك أكبر بثلاث مرات يستغرق زمنًا أطول بتسع مرات، أي \qty{90}{h}.
\end{solution}

\begin{solution}{exo:b3:inorganic-materials:2}
الحلمأة: $\ce{Si(OC2H5)4 + 4H2O -> Si(OH)4 + 4C2H5OH}$.\\
التكاثف: $\ce{2Si(OH)4 -> Si2O(OH)6 + H2O}$.\\
التفاعل الإجمالي: $\ce{Si(OC2H5)4 + 2H2O -> SiO2 + 4C2H5OH}$.
\end{solution}

\begin{solution}{exo:b3:inorganic-materials:3}
المكوِّنات: $\ce{SiO2}$، $\ce{B2O3}$. والمعدِّلات: $\ce{Na2O}$، $\ce{CaO}$. يكسر كل أيون أكسيد من $\ce{Na2O}$ جسرًا Si--O--Si
إلى طرفين \babelsublr{Si--O$^-$}، يقترن كل منهما بأيون $\ce{Na+}$: فتنقطع الشبكة، ويسيل المصهور
بسهولة أكبر وينصهر عند درجة أدنى.
\end{solution}

\begin{solution}{exo:b3:inorganic-materials:4}
اثنا عشر خماسيًا (كما في كل \omterm{def:b3:inorganic-materials:carbon}{فوليرين}). $V = 70$، $E = 3V/2 = 105$ رابطة، $F = 2 - V + E = 37$
وجهًا، ومنه $37 - 12 = 25$ سداسيًا.
\end{solution}

\begin{solution}{exo:b3:inorganic-materials:5}
$\sqrt2\,(r_B + r_O) = 1.4142 \times \qty{200.5}{pm} = \qty{283.5}{pm}$. $\ce{SrTiO3}$: $284/283.5 = 1.00$، البيروفسكيت المكعب المثالي.
$\ce{BaTiO3}$: $301/283.5 = 1.06$، التيتانيوم أصغر من ثماني أوجهه، فينزاح عن المركز: \omterm{def:b3:inorganic-materials:ferroelectric}{مادة كهروحديدية}.
$\ce{CaTiO3}$: $274/283.5 = 0.97$، الكالسيوم أصغر من تجويفه: ثمانيات أوجه مائلة، وتناظر
أدنى.
\end{solution}

\begin{solution}{exo:b3:inorganic-materials:6}
Si/Al $= 106/86 = 1.23$. $M = 86 \times 23.0 + 86 \times 27.0 + 106 \times 28.1 + 384 \times 16.0 = \qty{13422.6}{g/mol}$؛ و43 أيون $\ce{Ca^2+}$ لكل
صيغة: $43/\qty{13422.6}{g/mol} = \qty{3.20}{mmol/g}$.
\end{solution}

\begin{solution}{exo:b3:inorganic-materials:7}
المسافة بين أقرب جارين $d = a/\sqrt2 = \qty{288.4}{pm}$. وقطر عنقود من $n$ قشرة يساوي $(2n + 1)d$:
$2n + 1 \approx 5/0.2884 = 17.3$، ومنه $n = 8$ (\qty{4.9}{nm}). $N(8) = 2057$ ذرة، منها $10 \times 64 + 2 = 642$ على السطح:
31\,\%.
\end{solution}

\begin{solution}{exo:b3:inorganic-materials:8}
$h^2/(8\mu R^2)$ مع $\mu = \qty{9.109e-32}{kg}$: \qty{0.94}{eV} عند \qty{2.0}{nm} و\qty{0.42}{eV} عند \qty{3.0}{nm}. والإصدار عند
$1.74 + 0.94 = \qty{2.68}{eV}$، أي \qty{463}{nm} (أزرق)، وعند $1.74 + 0.42 = \qty{2.16}{eV}$، أي \qty{574}{nm} (أصفر): فالنقطة الأصغر
تصدر الضوء الأزرق.
\end{solution}

\begin{solution}{exo:b3:inorganic-materials:9}
$(10, 10)$: كرسي، $n - m = 0$، فلزي. $(10, 0)$: متعرج، 10 ليس مضاعفًا للعدد 3، \omterm{def:b3:solid-state:classes}{شبه موصل}.
$(9, 0)$: متعرج، فلزي. $(12, 8)$: لا هذا ولا ذاك (كيرالي)، $n - m = 4$، \omterm{def:b3:solid-state:classes}{شبه موصل}.
\end{solution}

\begin{solution}{exo:b3:inorganic-materials:10}
$M = 4 \times 65.4 + 13 \times 16.0 + 24 \times 12.0 + 12 \times 1.0 = \qty{769.6}{g/mol}$؛ $a^3 = (\qty{25.82e-8}{cm})^3 = \qty{1.721e-20}{cm^3}$؛
$\rho = 8 \times 769.6/(\num{6.022e23} \times \num{1.721e-20}) = \qty{0.594}{g.cm^{-3}}$. ومساحة الملعب \qty{7140}{m^2}: فالغرام الواحد يعادل
$3000/7140 = 0.42$ ملعب، وملعقة صغيرة من بضعة غرامات أكثر من ملعب كامل.
\end{solution}

\begin{solution}{exo:b3:inorganic-materials:11}
$\dfrac{10n^2 + 2}{(10n^3 + 15n^2 + 11n + 3)/3} = \dfrac{30n^2 + 6}{10n^3 + 15n^2 + 11n + 3} \to \dfrac{30n^2}{10n^3} = \dfrac3n$. ومن أجل $n = 8$: النسبة الدقيقة
$642/2057 = 0.31$، مقابل $3/8 = 0.375$؛ فالتقريب يبالغ في تقدير العناقيد الصغيرة،
التي لا تُهمل قشورها الداخلية.
\end{solution}

\begin{solution}{exo:b3:inorganic-materials:12}
$2E = 3V$، $2E = 5p + 6h + 7s$، $F = p + h + s$. وأويلر: $V - E + F = 2$، مع $V = 2E/3$: $F - E/3 = 2$، ومنه
$p + h + s - (5p + 6h + 7s)/6 = 2$، أي $(p - s)/6 = 2$: $p - s = 12$. فكل سباعي يحتاج إلى خماسي
إضافي.
\end{solution}

\begin{solution}{pb:b3:inorganic-materials:1}
\textbf{1.} $\mathrm{Na_{12}Al_{12}Si_{12}O_{48}}$: $12 \times 23.0 + 12 \times 27.0 + 12 \times 28.1 + 48 \times 16.0 = \qty{1705.2}{g/mol}$.

\textbf{2.} $1705.2 + 27 \times 18.0 = \qty{2191.2}{g/mol}$.

\textbf{3.} Si/Al $= 12/12 = 1$.

\textbf{4.} لا روابط Al--O--Al؛ ومع Si/Al $= 1$ يكون لكل Si أربعة جيران Al ولكل
Al أربعة جيران Si: تناوب صارم.

\textbf{5.} $-12$، شحنة واحدة لكل ألومنيوم.

\textbf{6.} $486.0/2191.2 = 22.2$\,\%.

\textbf{7.} $\ce{2Na+}(\text{الزيوليت}) + \ce{Ca^2+}(\text{aq}) \rightleftharpoons \ce{Ca^2+}(\text{الزيوليت}) + \ce{2Na+}(\text{aq})$.

\textbf{8.} ستة (اثنتا عشرة شحنة).

\textbf{9.} $6/\qty{1705.2}{g/mol} = \qty{3.52}{mmol/g}$.

\textbf{10.} $6/\qty{2191.2}{g/mol} = \qty{2.74}{mmol/g}$.

\textbf{11.} $\qty{3.52}{mmol/g} \times \qty{100.1}{g/mol} = \qty{352}{mg}$ من $\ce{CaCO3}$ لكل غرام.

\textbf{12.} $\qty{2.0}{mmol/L} \times \qty{15}{L} = \qty{30}{mmol}$.

\textbf{13.} $\qty{30}{mmol}/\qty{2.74}{mmol/g} = \qty{11}{g}$.

\textbf{14.} $\qty{10.9}{g}/0.60 = \qty{18}{g}$.

\textbf{15.} $\qty{60}{mmol}$ من $\ce{Na+}$، أي $\qty{0.060}{mol} \times \qty{23.0}{g/mol} = \qty{1.4}{g}$.

\textbf{16.} الفوسفات في المياه المستعملة يغذي الطحالب في الأنهار والبحيرات؛ 
وتكاثرها ثم تحللها يستنفدان الأكسجين المذاب (الإثراء الغذائي).

\textbf{17.} يمسك الأيون $\ce{Mg^2+}$ الأصغر جزيئات مائه بقوة أكبر؛ ويجب أن يتخلص منها
ليعبر النافذة ويبلغ المواقع، وهذا بطيء عند درجات حرارة الغسيل.

\textbf{18.} قطر الأيون العاري نحو \qty{0.20}{nm}، أي نصف النافذة.

\textbf{19.} الأيون المميَّه أكبر من النافذة: فيجب أن يتخلص من جزء من
غلافه المائي ليعبر، وتحل ذرات أكسجين الهيكل محل
الماء.

\textbf{20.} الهيكل مشحون سالبًا ولا يبادل إلا الكاتيونات؛ والأنيون الكبير
ذو السلسلة الطويلة يُصد ولا يستطيع عبور النافذة، فيبقى
في الماء، حيث يؤدي عمله.

\textbf{21.} \omterm{def:b3:inorganic-materials:zeolite}{الزيوليت} المنزوع الماء يأخذ الماء بقوة إلى أقفاصه؛ ولا يدخل إلا
الجزيئات الأصغر من النافذة، فهو يفصل الجزيئات بحسب حجمها:
منخل على المقياس الجزيئي.

\textbf{22.} تجلس أيونات $\ce{K+}$ الأكبر في النوافذ وتضيّقها: فلا تُقبل إلا الجزيئات
الأصغر، مثل الماء.

\textbf{23.} المتماكب \textit{para} النحيل الخطي يتسع في القنوات وينتشر بسرعة؛ أما
المتماكبان \textit{ortho} و\textit{meta} الأعرض فينتشران ببطء، ويبقيان ويتماكبان.

\textbf{24.} يستطيع \omterm{def:b3:inorganic-materials:zeolite}{الزيوليت} A اللامائي، نظريًا، أن يأخذ \qty{3.52}{mmol} من $\ce{Ca^2+}$ لكل غرام.
\end{solution}
